\chapter{Goals and Non-Goals}
\label{ch:goals}

\section{Goals}

The following goals are the binding commitments of the design. Where possible, each goal is stated
as a verifiable criterion rather than an aspiration.

\begin{enumerate}[label=\textbf{G\arabic*.}]

  \item \textbf{Preserve and formalize the governance layer.}
    Every capability that made v1 operationally useful---command allowlists and denylists, environment
    variable blocking, output redaction with regex patterns, and per-step role-based approval
    gates---must be present with identical or stronger semantics. Governance is gert's core
    differentiator~\cite{pagerduty-automation,aws-ssm-runbooks}: approval gates
    integrated directly into step definitions. must additionally support RBAC scoping on runbook
    execution: which identities may run which runbooks, and under what conditions.

  \item \textbf{Preserve the append-only, crash-safe, tamper-evident trace.}
    The execution trace must be an append-only JSONL file. Each line is a self-describing JSON
    object with a unique event ID (UUID), ISO~8601 UTC timestamp~\cite{rfc3339}, actor identity, step ID, and
    structured payload. The file must be crash-safe (partial writes must not corrupt prior entries)
    and must support optional cryptographic signing to satisfy SOC~2 and ISO~27001~(A.12.4) audit
    requirements~\cite{soc2-aicpa,iso27001}. This is a launch blocker: no trace, no compliance.

  \item \textbf{Preserve evidence capture, deterministic replay, TUI, VS Code extension, and
    \texttt{gert test}.}
    The following v1 capabilities must be functionally preserved now includes: evidence capture with text,
    checklists, and SHA256-hashed attachments; scenario-based deterministic replay (\texttt{--mode
    replay}); the Bubble~Tea terminal UI~\cite{bubbletea-framework}; the VS~Code extension powered by the JSON-RPC server~\cite{jsonrpc-spec,vscode-extension-api}; and
    the \texttt{gert test} scenario runner with its assertion engine. These features represent
    validated operational value. Their absence, is a regression, not a simplification.

  \item \textbf{Establish stable, versioned extension contracts.}
    Tool definitions (\texttt{.tool.yaml}), input providers (\texttt{.provider.yaml}), and
    out-of-process extensions must bind to named, versioned contracts. A host version bump that breaks
    an extension contract requires a major version increment and a migration path. The capability
    taxonomy---what an extension may register, observe, and invoke---must be a published, enumerated
    list, not an implicit set of permissions inferred from code. This pattern is established by
    VS~Code's extension API~\cite{vscode-extension-api} and Kubernetes admission controller design~\cite{kubernetes-docs}.

  \item \textbf{Introduce saga and compensation for multi-step rollback.}
    must support explicit compensation registration: when a forward step completes, the runbook may
    declare a compensating action to execute in reverse order on failure. This is the industry-standard
    saga pattern, implemented in production by Temporal~\cite{temporal-docs} and AWS Step Functions~\cite{aws-step-functions}. Without compensation,
    partially-executed runbooks that modify infrastructure leave systems in undefined intermediate
    states.

  \item \textbf{Integrate OpenTelemetry for distributed tracing.}
    Every runbook execution must emit OpenTelemetry spans~\cite{otel-spec}: one root span per run, one child span per
    step, with W3C Trace Context~\cite{w3c-trace-context} (\texttt{traceparent} / \texttt{tracestate}) propagated to tool
    invocations. Correlation IDs must appear in both the JSONL trace and the OTel span attributes, so
    that run traces can be joined with infrastructure observability data in Jaeger~\cite{jaeger-docs}, Zipkin~\cite{zipkin-docs}, or any
    OTLP-compatible backend. This is the current industry standard for operational visibility.

  \item \textbf{Decouple adapters from core via a typed event API.}
    The TUI, VS~Code extension, and any future web adapter must attach to the core exclusively via a
    typed event subscription interface. No adapter may import internal runtime packages or call runtime
    methods directly. The event schema must be versioned: adapters declare which schema version they
    consume, and the host rejects incompatible connections. This constraint is necessary for gert to
    support concurrent adapters (e.g., TUI and VS~Code running simultaneously against a single run).

  \item \textbf{Establish migration compatibility for v1 runbooks.}
    The command \texttt{gert migrate} must produce a valid \texttt{runbook/v1} document
    from any well-formed \texttt{runbook/v0} or \texttt{runbook/v1} input with no manual fixup
    required for standard runbooks. The runtime should additionally support a compatibility
    execution mode that accepts v1 runbooks without prior migration, so that existing users are not
    blocked from adopting the schema incrementally.

\end{enumerate}

\section{Non-Goals}

The following are explicitly outside the scope of gert. These boundaries are feature-scoped, not
process-scoped, and are intended to prevent scope creep during implementation.

\begin{enumerate}[label=\textbf{NG\arabic*.}]

  \item \textbf{Gert does not provide a managed cloud execution service.}
    There is no gert-as-a-service, no central execution registry, and no SaaS offering. Gert is a
    local-first binary. Cloud scheduling, multi-tenant isolation, and remote execution orchestration
    are concerns for a hypothetical v3 or a separate product.

  \item \textbf{Gert does not implement a general-purpose workflow orchestrator.}
    Gert is not Temporal, Argo, or Prefect~\cite{temporal-docs,argo-workflows,prefect-docs}. It does not support distributed workers, fan-out
    parallelism across remote nodes, or arbitrary DAG execution. Parallel execution within a single
    local run is a stretch goal; cross-host parallelism is explicitly out of scope.

  \item \textbf{Gert does not ship a built-in policy server.}
    This design defines governance semantics and may integrate policy-as-code hooks~\cite{opa-docs}, it does not
    bundle or operate an OPA server. Policy evaluation happens at the host process level. Central
    policy distribution, policy bundles over HTTP, and multi-host policy synchronization are out of
    scope.

  \item \textbf{Gert does not define an extension marketplace or signing authority.}
    The runtime will define an extension signing \textit{mechanism} (a verifiable public key model), but it
    will not operate a registry, curate extensions, or act as a certificate authority. Extension
    distribution is the responsibility of the extension author.

  \item \textbf{Gert does not guarantee binary stability for internal packages.}
    The internal Go package structure, private interfaces, and runtime internals are not subject to
    backward compatibility guarantees. Only the documented public contracts---the runbook schema,
    the tool schema, the JSON-RPC API surface, the extension protocol, and the JSONL trace
    schema---carry compatibility commitments.

\end{enumerate}
